\section{Measurable labels and finite blocks}\label{sec:labels}

We retain the saturated array object of Section~\ref{sec:arrays}.
Its local conditional laws are $\eta_u$, its parameter-tuple law is $\pi$,
and its conditional state-tuple laws are $\lambda_p$.
Proposition~\ref{prop:common-exact-lines} describes intersections of active
spaces in a square. We must turn that intersection statement into a
description on each individual fibre. First we show that an entire
active space is spanned by exact lines. We then enumerate those lines
measurably over tuple parameters and prove that the family depends only
on the local parameter. Finally, we partition the lines into finite
sets on which time has a fixed permutation rule. One such moving set
will retain a nonzero coefficient of the original mixing witness,
providing the input for the cancellation argument in
Section~\ref{sec:cancellation}.

Write $H_u=L^2(\eta_u)$. For a tuple $p=(p_1,\ldots,p_k)$ and a slot
$i\in I$, let $\Pi_i(p)$ be the active projection in $H_{p_i}$ for
$\lambda_p$, and let $C_i(p)$ be the closed span of its exact lines.
At this point $C_i(p)$ is defined separately for each tuple; no
measurability of its line family is assumed.
We use $p_a[j]$ for the entry of row tuple $p_a$ in column $j$, so that
the row index on $p_a$ is distinguished from a coordinate of a single
tuple. Let $\vartheta$ denote the law of the complete parameter square
$E=(u_{a,j})_{a,j\in I}$.

\subsection{From intersections to entire active spaces}

The fields $H_u$ can be coordinated measurably using a countable algebra
of bounded simple functions on the standard state space and
Gram--Schmidt orthogonalization. Zero vectors are omitted, or retained
as zero padding in a fixed coordinate space $\ell^2$. These local
coordinates depend on $u$ alone. In particular, two tuple parameters
with the same entry $u$ refer to the same coordinated Hilbert space.
Conditional-expectation contractions of countably many product tests
give a measurable spanning family for each active space. Increasing
finite-dimensional projections onto these spans show that
$p\mapsto\Pi_i(p)$ is a measurable operator field.

\begin{lemma}[Active-space exactness]\label{lem:active-exact}
For $\pi$-almost every $p$ and every $i\in I$,
\begin{equation}\label{eq:active-exact}
 \Ran\Pi_i(p)=C_i(p).
\end{equation}
\end{lemma}

\begin{proof}
Fix $i$. By Lemma~\ref{lem:parameter-geometry}, the row and column
tuples meeting at cell $(i,i)$ are conditionally independent with law
\[
 \pi_i^u:=\pi(\dd p\mid p_i=u)
\]
given their common entry $u$. Proposition~\ref{prop:common-exact-lines}
therefore gives, for almost every such pair $(p,q)$,
\begin{equation}\label{eq:random-projection-intersection}
 \Pi_i(p)\Pi_i(q)=\Pi_i(q)\Pi_i(p),\qquad
 \Ran\Pi_i(p)\cap\Ran\Pi_i(q)\subseteq C_i(p).
\end{equation}

Here the almost-everywhere assertion does not presuppose a measurable
description of $C_i$. Indeed, choose a measurable conull set of
parameter squares on which the conditional kernel hypotheses of
Proposition~\ref{prop:common-exact-lines} hold on countable determining
classes. That Proposition gives its pointwise conclusion on every
square in this set. Disintegrate this set over the two tuples meeting
at $(i,i)$. On a measurable conull set of pairs the conditional measure
assigns probability one to admissible squares, and is concentrated on
the prescribed pair. Each such pair has an admissible completion.
The conclusions in \eqref{eq:random-projection-intersection} depend
only on the two tuple laws, so they hold for that pair. No measurable
choice of a completion or of an exact line is required. Disintegration
over $u$ then gives conull sections for the product law
$\pi_i^u\otimes\pi_i^u$.

For a typical $u$ define the positive contraction on $H_u$
\begin{equation}\label{eq:mean-active-projection}
 B_{i,u}=\int \Pi_i(q)\,\pi_i^u(\dd q),
 \qquad Q_{i,u}=\text{the projection onto }\Ker B_{i,u}.
\end{equation}
The weak operator integral is well defined, and its matrix entries are
measurable in $u$. Since $0\le B_{i,u}\le\id$,
$Q_{i,u}$ is the strong limit of $(\id-B_{i,u})^n$ and is measurable too.
In the fixed ambient $\ell^2$ coordinates, use the projection of each
standard basis vector $e_j$ into $H_u$ and extend the operators by zero
on $H_u^\perp$. Then
\[
 \int\norm{\Pi_i(p)Q_{i,u}e_j}^2\,\pi_i^u(\dd p)
 =\ip{Q_{i,u}e_j}{B_{i,u}Q_{i,u}e_j}=0.
\]
After a countable intersection, $\Pi_i(p)Q_{i,u}=0$ for almost every
$p$. Taking adjoints gives
\begin{equation}\label{eq:active-support}
 \Ran\Pi_i(p)\subseteq(\Ker B_{i,u})^\perp.
\end{equation}
This range containment is an operator assertion on a measurable conull
set: it is tested by $Q_{i,u}\Pi_i(p)e_j=0$ for all $j$.

Fix $p$ satisfying \eqref{eq:active-support} and having a conull
$q$-section in \eqref{eq:random-projection-intersection}.
If $v\in\Ran\Pi_i(p)$ is orthogonal to $C_i(p)$, commutation implies
that $\Pi_i(q)v$ belongs to the intersection in
\eqref{eq:random-projection-intersection}. Consequently
\[
 \norm{\Pi_i(q)v}^2=\ip{v}{\Pi_i(q)v}=0
 \quad\text{for }\pi_i^u\text{-almost every }q.
\]
Thus $\ip{v}{B_{i,u}v}=0$, and $v\in\Ker B_{i,u}$.
Equation~\eqref{eq:active-support} forces $v=0$.
The conull $q$-section used here gives whole-subspace statements, so it
works for every such $v$; a measurable selection of $v$ is unnecessary.
Every exact line is active, by its defining product relation. This
proves \eqref{eq:active-exact}. There are finitely many slots $i$.
\end{proof}

\subsection{A measurable group of lines over each local parameter}

A measurable countable family over $U$ will mean a standard measurable
space over $U$ whose fibres admit countably many measurable partial
sections enumerating them without repetition. In this description,
measurability of a fibrewise operation means measurability of its
indices in these enumerations. A line is represented by a unit vector;
for the exact lines below such a vector has modulus one almost surely.
Different scalar phases represent the same line.

\begin{proposition}[Measurable local labels]\label{prop:measurable-labels}
There is a measurable countable family of abelian groups
$u\mapsto\mathcal L(u)$, defined on a $\nu$-conull set, with the
following properties.
\begin{enumerate}
\item Each $\ell\in\mathcal L(u)$ labels a line in $H_u$ with a
measurable modulus-one representative $\chi_{u,\ell}$. The identity
label $0$ has representative $1$. Distinct labels are orthogonal,
every nonidentity representative has mean zero, and
\[
 \chi_{u,\ell}\chi_{u,\ell'}
 \in\T\chi_{u,\ell+\ell'}.
\]
\item For $\pi$-almost every $p$, the exact lines at slot $i$ of
$\lambda_p$ are precisely $\mathcal L(p_i)$, for every $i$, and they
span its active space. Unique exact matching gives measurable group
isomorphisms
\[
 \Phi_{st}(p):\mathcal L(p_s)\longrightarrow\mathcal L(p_t)
 \qquad(s,t\in I),
\]
where $\Phi_{ss}(p)=\id$.
\item Local time gives measurable group isomorphisms
\[
 \tau_u:\mathcal L(u)\longrightarrow\mathcal L(T_Uu),
 \qquad [\chi]\longmapsto[\chi\circ T_Y^{-1}].
\]
They respect exact matching. All local time assertions may be imposed
on an invariant conull set.
\item For $\vartheta$-almost every parameter square, prescribing a
label at any one cell gives a unique array of labels matched along
every row and every column.
\end{enumerate}
\end{proposition}

\begin{proof}
We first enumerate the exact lines over tuple parameters, then show
that these enumerations describe a family depending only on a local
entry. The distinction is necessary: Lemma~\ref{lem:active-exact}
alone gives no such independence from the other parameters.

\paragraph{Enumerating the tuple-dependent lines.}
Fix $i$ and another slot $j\ne i$. Let
\[
 \Delta_i(p):H_{p_i}\longrightarrow
       \bigotimes_{t\ne i}H_{p_t}
\]
be conditional expectation onto all the other coordinates under
$\lambda_p$. Their joint marginal is product, so the displayed target
is their $L^2$ space. This is a measurable contraction field. It is the
adjoint of conditional expectation in the other direction, and hence
vanishes on the inactive complement. For a unit representative of an
exact line, $\Delta_i(p)$ gives, up to a scalar phase, the unit tensor
of the conjugated matched representatives in the other slots.

Choose dense unit-ball test sections $g_n(u)$ in each $H_u$. If
$\theta_g h=g\ip{g}{h}$, set
\begin{equation}\label{eq:line-detection-operators}
 A_n(p)=\Delta_i(p)^*
   \bigl(\theta_{g_n(p_j)}\otimes\id_{t\ne i,j}\bigr)
   \Delta_i(p).
\end{equation}
All their matrix entries are measurable before any exact-line
enumeration. Pointwise, Lemma~\ref{lem:active-exact} and
Lemma~\ref{lem:exact-lines} give an orthogonal exact-line basis of the
active space. Write $w_{j,\ell}$ for the unit partner appearing at $j$
in $\Delta_i(p)\chi_\ell$. The operators in
\eqref{eq:line-detection-operators} are diagonal on this basis, with
eigenvalues
\[
 a_n(\ell)=\abs{\ip{g_n(p_j)}{w_{j,\ell}}}^2.
\]
Indeed, for two distinct labels the partners are orthogonal in every
other slot. Since $k\ge3$, at least one slot outside $i,j$ is untouched
by the insertion, and its inner product kills the off-diagonal entry.
The partners at $j$ are orthonormal, so Bessel's inequality gives
\begin{equation}\label{eq:line-operator-trace}
 \tr A_n(p)=\sum_\ell a_n(\ell)\le1.
\end{equation}
Dense tests detect each line and distinguish any two different lines:
the corresponding rank-one projections at $j$ are different.

Here is a measurable spectral construction which does not assume
that the basis just used pointwise is already enumerated measurably.
Choose independent continuously distributed weights
$w_n\in(0,2^{-n})$ and form
\[
 A(p,w)=\sum_{n\ge1}w_n A_n(p).
\]
The series converges in operator norm and is jointly measurable.
It is positive and has trace at most $\sum_n w_n$ on the conull set
where the preceding arguments apply. Trace finiteness is itself
measurable, being the finiteness of the sum of nonnegative diagonal
entries in the fixed coordinate space. Set the operator to zero on
its null exceptional set, so that the ensuing recursion always acts
on a positive compact operator.

For such an operator $A$, its norm is measurable from a countable
dense set of unit vectors. If $a=\norm A>0$, the strong limit
\[
 Q=\lim_{m\to\infty}(A/a)^m
\]
is its top-eigenvalue projection. Its matrix entries and rank are
measurable. Remove $aQ$ and repeat, taking zero projections if the
remainder is zero. Compactness ensures that this lists every distinct
positive eigenvalue and its finite-dimensional eigenspace. Thus a
repeated positive eigenvalue is detected by the measurable event that
some nonzero projection in the recursion has rank at least two.

For a fixed $p$, two distinct exact lines have different coefficient
sequences $(a_n(\ell))_n$. Conditioning on all but a weight where they
differ shows that their weighted eigenvalues coincide with probability
zero. There are only countably many lines in the separable Hilbert
space. Fubini applied to the measurable collision event therefore
chooses one deterministic sequence of weights for which the positive
spectrum is simple for almost every $p$. Every active line has a
strictly positive eigenvalue, because some test detects it and all
weights are positive. The zero eigenspace is precisely the inactive
complement; eigenvalues tending to zero add no kernel vectors. The
positive spectral projections therefore enumerate exactly the exact
lines. Zero padding of the Hilbert coordinates supplies no positive
eigenvectors and hence no extra lines.

Choose a unit vector in each one-dimensional projection by projecting
the first coordinate vector with nonzero image and normalizing.
These measurable Hilbert sections have jointly measurable function
representatives. One direct verification is to approximate a section
in each fibre within $2^{-m}$ by the first suitable member of a
countable family of rational combinations of simple sections.
The choices are measurable. The successive differences are summable
in the integrated $L^2$ norm, and hence in $L^1$, giving a measurable
almost-everywhere limit that represents the section. Each unit vector
on an exact line has modulus one, so its representative does also.
For a listed tuple of such representatives, their product $Z$ is
constant under $\lambda_p$ exactly when
\begin{equation}\label{eq:exact-tuple-test}
 \abs{\int Z\dd\lambda_p}=1.
\end{equation}
This measurable scalar test enumerates the unique matched tuples and
is unchanged by all choices of scalar phases.

\paragraph{Forcing all column memberships.}
Return to a typical parameter square. Fix a matched tuple along row
$a$, with unit representatives $v_j$ in its cells. At cell $(a,j)$ let
$Q_j$ denote the column-active projection $\Pi_a(q_j)$.
Proposition~\ref{prop:common-exact-lines} and orthogonality of distinct
row exact lines show that
\begin{equation}\label{eq:membership-indicator}
 Q_jv_j=\varepsilon_jv_j,\qquad \varepsilon_j\in\{0,1\}.
\end{equation}
Here $\varepsilon_j=1$ precisely when this row line is also an exact
line of the column law.
The row Markov operator between two cells, given the remaining cells,
intertwines their column-active projections. The exact relation maps
one chosen representative to the conjugate partner times a
modulus-one scalar depending on the remaining cells. Since the
projections commute with conjugation, this intertwining and
\eqref{eq:membership-indicator} imply
$\varepsilon_j=\varepsilon_t$ for every $j,t$.

Choose the row tuple measurably from the enumeration just obtained.
Given $p_a$, each $\varepsilon_j$ depends only on $q_j$ and this fixed
tuple. Lemma~\ref{lem:parameter-geometry} makes these indicators
pairwise conditionally independent. Their common conditional
probability $r(p_a)$ therefore satisfies $r(p_a)=r(p_a)^2$:
indeed, writing $r_j(p_a)=\E[\varepsilon_j\mid p_a]$, the almost-sure
equality of the indicators makes all $r_j$ equal, while conditional
independence gives $r_j=r_jr_t$ for $j\ne t$.
At the matched column $j=a$, however,
\[
 r(p_a)=\int\norm{\Pi_a(q)v_a}^2\,
             \pi(\dd q\mid q_a=p_a[a])
       =\ip{v_a}{B_{a,p_a[a]}v_a}>0.
\]
The strict inequality follows from \eqref{eq:active-support}, since
$v_a$ is a nonzero active vector. Thus $r(p_a)=1$.
Countably many row labels and finitely many cells make the resulting
inclusions simultaneous: every row exact line is a column exact line
at its cell, almost surely in the parameter square.

\paragraph{Descent to the local parameter.}
At $(i,i)$ two conditionally iid families of exact lines over the
shared parameter $u$ satisfy inclusion almost surely. Swapping the
two copies gives reverse inclusion, so the families agree.
To descend this agreement measurably, for each local dense test $g_n(u)$
encode a family $\mathcal S$ of orthogonal unit lines by
\begin{equation}\label{eq:line-family-encoding}
 K_n(\mathcal S)=
   \sum_{v\in\mathcal S}\abs{\ip{g_n(u)}v}^2\theta_v.
\end{equation}
The sum is independent of phases and order, has trace at most one,
and is measurable for our tuple-dependent enumeration. For every
matrix entry, the two conditional iid encodings are equal almost
surely. Their conditional squared difference is twice the conditional
variance, so that variance is zero. Each entry equals its conditional
expectation given $u$. The encodings consequently give measurable
operator fields depending only on $u$.
The family of encodings detects and distinguishes its orthogonal
lines, so the same positive-weight spectral construction recovers a
measurable family $\mathcal L_i(u)$ agreeing with the slot-$i$ exact
lines for almost every tuple with that entry.

At a nonmatched cell $(a,j)$ the already proved inclusion now reads
$\mathcal L_j(u_{a,j})\subseteq\mathcal L_a(u_{a,j})$.
The marginal of $u_{a,j}$ is $\nu$. Applying this to every ordered
pair $a,j$ yields one common family $\mathcal L(u)$ almost everywhere.
We have now removed the dependence on the other entries of a tuple:
the remaining construction can be made on the local parameter space.

By Lemma~\ref{lem:exact-lines}, multiplication of representatives
modulo scalars makes each family an abelian group. Orthogonality and
the zero-mean assertion for nonidentity labels follow from the same
Lemma. Equality of unit lines is tested by $\abs{\ip v w}=1$.
Testing the product of two representatives against the countable list
therefore makes multiplication measurable; conjugation gives the
inverse. The constant line is measurably identifiable and we choose
its representative to be $1$. Equation~\eqref{eq:exact-tuple-test}
and uniqueness give the measurable group isomorphisms $\Phi_{st}(p)$.

The equivariant versions of $\eta_u$ and $\lambda_p$ show that
composition with $T_Y^{-1}$ carries exact tuples to exact tuples.
It consequently induces $\tau_u$ and respects all matchings.
Equality tests against the measurable line lists show that this map
and its inverse are measurable. Intersecting local conull sets over
all integer translates makes the local assertions time invariant.
Tuple and square assertions are imposed separately on their own
conull sets, also intersected over integer time translates; nothing
asserts validity for every tuple over the local conull set.

\paragraph{Consistency around the square.}
Start with a matched label tuple along a row and extend every one of
its labels along its column by exact matching. Choose representatives
for the resulting labels at all cells. The product down each column
is constant under the conditional square law, as is the product
along the initial row. Dividing the product of the column relations
by that row relation shows that the product of the remaining $k-1$
row products is constant. Those rows are independent as wholes by
Proposition~\ref{prop:square-locality}. Each row product has modulus
one. The absolute mean of their product is one, so each has absolute
mean one and is itself constant. Thus every row is matched.
Starting from a single cell, its row and then all columns are uniquely
determined; the preceding argument supplies consistency along the
other rows. Countability of the label lists makes this simultaneous
for every initial label. This proves the final assertion.
\end{proof}

We now have a countable group of lines over each local parameter,
with matchings consistent across the square. It remains to find a
finite moving set of labels that retains the nonzero level-$1$
correlation.

\subsection{Counting labels and their transition operators}

We next use the parameter-square geometry to compare labels over
different local parameters. Counting measure is the appropriate
measure on each group: exact matching is a bijection, even when the
group is infinite.
Set
\[
 Z=\{(u,\ell):\ell\in\mathcal L(u)\},\qquad
 \int_Z F\dd\xi=
 \int_U\sum_{\ell\in\mathcal L(u)}F(u,\ell)\dd\nu(u)
 \quad(F\ge0).
\]
The measurable enumeration makes $\xi$ sigma-finite. The measurable
map $\tau(u,\ell)=(T_Uu,\tau_u\ell)$ preserves $\xi$, because time
preserves $\nu$ and bijects the label fibres.

\begin{lemma}[Counting lift]\label{lem:counting-lift}
There are sigma-finite measures $\widetilde\pi$ on $Z^I$ and
$\widetilde\vartheta$ on $Z^{I\times I}$, obtained by counting matched
label tuples and arrays over $\pi$ and $\vartheta$, respectively.
Every coordinate marginal is $\xi$, and each row and column marginal
of $\widetilde\vartheta$ is $\widetilde\pi$.
For $s\ne t$ the coordinate transition in $\widetilde\pi$ is the
probability kernel
\begin{equation}\label{eq:label-transition-kernel}
 K_{st}((u,\ell),A)=
 \int\mathbf1_A\bigl(p_t,\Phi_{st}(p)\ell\bigr)
                   \,\pi(\dd p\mid p_s=u).
\end{equation}
Its operator $M_{st}$ on $L^2(\xi)$ is a contraction, has adjoint
$M_{ts}$, and commutes with $V F=F\circ\tau$. For all $a\ne b$,
\begin{equation}\label{eq:label-square-identity}
 M_{st}=M_{ab}^*M_{st}M_{ab}.
\end{equation}
Given a labelled row in $\widetilde\vartheta$, any proper family of
labelled columns has the product of the column laws conditioned on
their respective shared labelled points.
\end{lemma}

\begin{proof}
For a nonnegative measurable function $F$ on $Z^I$, define
\begin{equation}\label{eq:tuple-counting-lift}
 \int F\dd\widetilde\pi
 =\int\sum_{\ell\in\mathcal L(p_s)}
 F\bigl((p_t,\Phi_{st}(p)\ell)_{t\in I}\bigr)\dd\pi(p).
\end{equation}
Unique exact matching makes this independent of the chosen source
slot $s$. Define the square measure in the same way, counting its
unique matched arrays from one chosen cell. The last assertion of
Proposition~\ref{prop:measurable-labels} makes this independent of
the chosen cell. Counting from a coordinate, row, or column gives
the claimed marginals. The preimages of the first $N$ labels of any
one coordinate have measure at most $N$ and exhaust the lifted
spaces. Thus the measures are sigma-finite.

Disintegrating \eqref{eq:tuple-counting-lift} first over $p_s=u$
gives \eqref{eq:label-transition-kernel}. For each prescribed source
label there is exactly one completion over every compatible tuple;
fixing that label introduces no weight into the parameter kernel.
This proves directly that the conditional kernels are probability
kernels even when the total lifted measure is infinite.
Similarly, given a labelled row of the square, its label has exactly
one extension over every compatible parameter square. Its conditional
parameter-square law is therefore the original law given that row's
parameter tuple. Lemma~\ref{lem:parameter-geometry} gives independent
proper columns; their labels are then determined separately by their
shared labels and their own parameter tuples. This proves the last
assertion of the Lemma.

Write $\widetilde\pi_{st}$ for the $(s,t)$ marginal.
Its two marginals are $\xi$, and its kernels in the two directions
are $K_{st}$ and $K_{ts}$. Jensen's inequality and the marginal
identities give $L^2$ contraction and $M_{st}^*=M_{ts}$.
Time equivariance of exact matching preserves both lifted measures
and gives $VM_{st}=M_{st}V$.

Finally, test the $(s,t)$ pair in row $b$ of the lifted square with
$f,g\in L^2(\xi)$. Conditional on labelled row $a$, its two column
transitions are independent applications of $M_{ab}$, since two
columns form a proper family when $k\ge3$. Hence
\begin{align*}
 \ip f{M_{st}g}
 &=\int\overline{f(z_{b,s})}\,g(z_{b,t})
                              \dd\widetilde\vartheta\\
 &=\int\overline{(M_{ab}f)(z_{a,s})}\,
                  (M_{ab}g)(z_{a,t})\dd\widetilde\vartheta\\
 &=\ip{M_{ab}f}{M_{st}M_{ab}g}
  =\ip f{M_{ab}^*M_{st}M_{ab}g}.
\end{align*}
These integrals are absolutely integrable by Cauchy--Schwarz with
the pair marginals $\xi$. Alternatively one may first use bounded
tests of finite-measure support and then pass in $L^2$.
The calculation proves \eqref{eq:label-square-identity} without
normalizing the sigma-finite measures.
\end{proof}

\subsection{Finite blocks and the time permutation}

The transition operators constrain how the countable groups can vary
with the parameter. We now extract finite pieces that have fixed sizes
and a common permutation under time. Their individual labels need not
follow a parameter-independent time rule. The construction starts with
the information transported deterministically by a transition operator.
Independence of the local parameters will make its possible values
countable; a positive $L^2$ test will make each corresponding set of
labels finite. The group matchings will then make these finite sizes
the same in almost every fibre.

\begin{lemma}[Finite label blocks]\label{lem:finite-blocks}
There are a countable index set $\mathcal B$, measurable partitions
\[
 \mathcal L(u)=\bigsqcup_{b\in\mathcal B}B_b(u)
 \quad\text{for }\nu\text{-almost every }u,
\]
and permutations $D_0,S_0$ of $\mathcal B$ with these properties.
Every $B_b(u)$ has a fixed finite positive cardinality $m_b$, independent
of $u$. For one fixed pair $s\ne t$, exact matching satisfies
\[
 \Phi_{st}(p)B_b(p_s)=B_{D_0b}(p_t)
 \quad(\pi\text{-almost every }p),
\]
and time satisfies
\[
 \tau_u B_b(u)=B_{S_0b}(T_Uu).
\]
Moreover, there are a countable abelian group $\mathcal L$, a
partition $\mathcal L=\bigsqcup_b B_b$ into sets of cardinalities $m_b$,
block-preserving group isomorphisms
$J_u:\mathcal L\to\mathcal L(u)$ for typical $u$, and an automorphism
$G$ of $\mathcal L$ such that
\begin{equation}\label{eq:abstract-block-time}
 G^r B_b=B_{S_0^{-r}b}\qquad(r\in\Z).
\end{equation}
The isomorphisms $J_u$ are only asserted to exist separately in each
fibre; they need not be chosen measurably or intertwine individual
time trajectories.
\end{lemma}

\begin{proof}
Fix $s\ne t$ and write $M=M_{st}$. Taking $(a,b)=(s,t)$ and $(t,s)$
in \eqref{eq:label-square-identity} gives
\begin{equation}\label{eq:partial-unitary-identities}
 M=M^*MM=MMM^*.
\end{equation}
Let $\mathcal H_1=\overline{\Ran M}$ and
$\mathcal H_2=\overline{\Ran M^*}$.
The first identity puts $\mathcal H_1$ in the fixed space of $M^*M$,
which lies in $\mathcal H_2$. The adjoint of the second identity puts
$\mathcal H_2$ in the fixed space of $MM^*$, which lies in
$\mathcal H_1$. They are therefore equal to a closed space
$\mathcal H$. On it $M^*M=MM^*=\id$, while
$M$ vanishes on $\mathcal H^\perp=\Ker M$.
Thus $M$ is unitary on $\mathcal H$, with inverse $M^*$ there.
For every $F\in\mathcal H$, equality in the conditional-expectation
norm gives
\begin{equation}\label{eq:deterministic-label-transition}
 F(z_t)=(MF)(z_s)
 \quad\widetilde\pi_{st}\text{-almost everywhere}.
\end{equation}
Indeed, the integral of the squared difference is
$\norm F_2^2-\norm{MF}_2^2=0$.

Choose a measurable enumeration $\ell_n(u)$ without repetition, and
set $h(u,\ell_n(u))=2^{-n}$ on its domain, starting at $n=1$.
Then $h>0$ and $\norm h_2^2\le\sum_{n\ge1}4^{-n}<\infty$.
The probability kernel sends $h$ to a strictly positive function
$v=Mh$ almost everywhere: the kernel avoids every $\xi$-null set
for almost every source point, by its marginal identity, and the
integral of a positive function under a probability law is positive.
Also $v\in\mathcal H$.

The time operator $V$ is unitary and commutes with $M$ in both
directions, so it preserves $\mathcal H$. Define the measurable
countable vector
\begin{equation}\label{eq:label-orbit-vector}
 \psi(z)=\bigl((M^nV^r v)(z)\bigr)_{(n,r)\in\Z^2}
          \in\C^{\Z^2}.
\end{equation}
Negative powers of $M$ are taken on $\mathcal H$.
All its coordinates have finite versions outside one null set, and
its $(0,0)$ coordinate is positive there. On $\C^{\Z^2}$ define
\[
 (D_0c)_{n,r}=c_{n+1,r},\qquad
 (S_0c)_{n,r}=c_{n,r+1}.
\]
These are commuting invertible coordinate reindexings.
Using \eqref{eq:deterministic-label-transition} for the countably
many coordinates gives
\begin{equation}\label{eq:psi-transports}
 \psi(z_t)=D_0\psi(z_s),\qquad
 \psi(\tau z)=S_0\psi(z)
\end{equation}
almost everywhere in the respective measures.

We claim that $\psi$ has countable essential image. Given a source
point $z=(u,\ell)$, the parameter of its $K_{st}$ target has law
$\nu$, since $p_s,p_t$ are independent under $\pi$ and fixing a
source label does not tilt that law. Consider the occurrence relation
\[
 R(z,u')\quad\Longleftrightarrow\quad
 D_0\psi(z)=\psi(u',\ell_n(u'))
 \text{ for some defined }\ell_n(u').
\]
It is measurable: existence is a countable union and equality of
vectors is a countable conjunction of coordinate equalities.
The first identity in \eqref{eq:psi-transports} implies
$\nu\{u':R(z,u')\}=1$ for $\xi$-almost every $z$.
Tonelli's theorem for the sigma-finite product $\xi\otimes\nu$
selects one $u'$ for which $R(z,u')$ holds for $\xi$-almost every
$z$. That target fibre has only countably many $\psi$-values.
Invertibility of $D_0$ proves the claim.

Let $\mathcal B$ be the positive-mass values in this countable
essential image. For $b\in\mathcal B$ put
\[
 E_b=\{z:\psi(z)=b\},\qquad
 B_b(u)=\{\ell:(u,\ell)\in E_b\}.
\]
Every $b$ has positive $(0,0)$ coordinate $b_{0,0}$, and
\begin{equation}\label{eq:finite-level-mass}
 0<\xi(E_b)\le \frac{\norm v_2^2}{|b_{0,0}|^2}<\infty.
\end{equation}
Equal marginals in the first identity of \eqref{eq:psi-transports}
and time invariance in the second show that $D_0$ and $S_0$ permute
$\mathcal B$, preserving each level's mass along their orbits.
After discarding a $\nu$-null set, all labels in every remaining
fibre belong to these levels, and each $B_b(u)$ is finite.
Countability of $\mathcal B$ makes these assertions simultaneous.

Counting the labels in \eqref{eq:psi-transports} shows that, for
$\pi$-almost every $p$, its matching isomorphism $\Phi_{st}(p)$
carries every block $b$ onto block $D_0b$.
Because each fibre's label list and $\mathcal B$ are countable,
intersecting the corresponding conull sets upgrades the transported
label identity to every label in every block for such a tuple.
Disintegrate the measurable conull set of such tuples over
$(p_s,p_t)$. Since its marginal is $\nu\otimes\nu$, there is a
measurable conull set of pairs $(u,w)$ admitting a tuple with this
property. Fubini chooses a common source $u_*$ for which almost every
target is in this set. For each such target $u$, choose one resulting
group isomorphism
$I_u:\mathcal L(u_*)\to\mathcal L(u)$.
Each $I_u$ takes block $b$ to block $D_0b$ simultaneously for every
$b$. The choices are used only pointwise.

Fix one typical target $u_0$, set $\mathcal L=\mathcal L(u_0)$,
and define $J_u=I_uI_{u_0}^{-1}$. These group isomorphisms preserve
every block index. Thus all $|B_b(u)|$ are constant almost everywhere.
Equation~\eqref{eq:finite-level-mass} makes each constant finite,
and its positive left side makes it nonzero. Let $B_b=B_b(u_0)$
and $m_b=|B_b|$.

The second identity of \eqref{eq:psi-transports} gives the asserted
time permutation of blocks. Restrict the local parameter set to
the intersection of all integer translates of the conull sets just
used and of the time-equivariance set. Choose $u$ in this invariant
conull set. The automorphism
\[
 A=J_{T_Uu}^{-1}\tau_uJ_u
\]
of $\mathcal L$ maps every $B_b$ to $B_{S_0b}$.
Take $G=A^{-1}$. Iteration gives
\eqref{eq:abstract-block-time}. This last step fixes one automorphism
realizing the permutation of blocks; it makes no assertion about
the maps inside a block at other parameters or later times.
\end{proof}

\subsection{A finite block carrying the witness}

All structural objects needed for the final argument are now
available. The next Proposition records the precise connection to
the nonzero correlation retained from the original process. Its
measurable quantities are sums over finite blocks, so the abstract
identifications in Lemma~\ref{lem:finite-blocks} will cause no
measurability issue.

\begin{proposition}[A block carrying the witness]\label{prop:block-witness}
Let the saturated array object lie over the zero-entropy finite-alphabet
mixing process $(X,\mu_X,T_X)$ supplied by
Proposition~\ref{prop:zero-entropy-witness}, with its centered
finite-valued one-site functions $f_1,\ldots,f_k$ and nonzero
level-$1$ product integral. There exist a function
$f\in\{f_1,\ldots,f_k\}$, a countable abelian group $\mathcal L$,
an automorphism $G$ of $\mathcal L$, a finite nonempty set
$B\subset\mathcal L$, and a number $c>0$ with the following properties.

For $\nu$-almost every $u$, the probability space $(Y,\eta_u)$
has factor map $x:Y\to X$ with $x_*\eta_u=\mu_X$, and admits
representatives $(\chi_\ell^u)_{\ell\in\mathcal L}$ such that
\begin{gather}
 \chi_0^u=1,\qquad |\chi_\ell^u|=1,\qquad
 \chi_\ell^u\chi_{\ell'}^u\in\T\chi_{\ell+\ell'}^u,
 \label{eq:block-character-laws}\\
 \int\chi_\ell^u\dd\eta_u=0\quad(\ell\ne0),\qquad
 \ip{\chi_\ell^u}{\chi_{\ell'}^u}_{L^2(\eta_u)}
       =\mathbf1_{\{\ell=\ell'\}}.
 \nonumber
\end{gather}
For every $r\ge0$ the function
\begin{equation}\label{eq:block-witness-sum}
 s_r(u)=\sum_{\ell\in B}
 \left|\int f(T_X^r x(y))\,
               \overline{\chi_{G^r\ell}^u(y)}\dd\eta_u(y)\right|
\end{equation}
has a measurable version and satisfies
\begin{equation}\label{eq:block-witness-integral}
 0\le s_r(u)\le |B|\norm f_\infty,
 \qquad \int_U s_r(u)\dd\nu(u)=c>0.
\end{equation}
The group identifications used to choose the representatives in
\eqref{eq:block-character-laws} need not be measurable in $u$.
Their finite sums in \eqref{eq:block-witness-sum} are intrinsic sums
over the measurable moving blocks. In particular, no time invariance
of an individual $\eta_u$ is asserted or required.
\end{proposition}

\begin{proof}
Condition the nonzero product integral at level $1$ on $p$.
For any chosen slot $i$, the conditional expectation of the product
of the other bounded inputs onto that slot belongs to its active
space. If $f_i(x)$ were orthogonal to that space for almost every
$p$, the product integral would vanish. In using complex inner
products here one can conjugate the other product; the active
space is conjugation invariant.
Lemma~\ref{lem:active-exact} and
Proposition~\ref{prop:measurable-labels} therefore give a choice
$f=f_i$ with a nonzero coefficient on some local label on a set of
positive $\nu$-measure. The identity-label coefficient is zero:
\[
 \int f(x(y))\dd\eta_u(y)=\int f\dd\mu_X=0.
\]
Thus nonidentity labels supply the nonzero coefficient.

For a block $b\in\mathcal B$ define the measurable function
\[
 c_b(u)=\sum_{\ell\in B_b(u)}
 \left|\int f(x(y))\overline{\chi_{u,\ell}(y)}\dd\eta_u(y)\right|.
\]
The countably many blocks cover all labels, so some fixed block $b$
has $\int c_b\dd\nu>0$. Its cardinality $m_b$ is fixed and finite,
and hence $c_b\le m_b\norm f_\infty$.
Set $B=B_b\subset\mathcal L$ using
Lemma~\ref{lem:finite-blocks}, and $c=\int c_b\dd\nu$.

For each typical $u$ choose the block-preserving identification $J_u$
from that Lemma and put
$\chi_\ell^u=\chi_{u,J_u\ell}$.
Proposition~\ref{prop:measurable-labels} gives
\eqref{eq:block-character-laws}.
By \eqref{eq:abstract-block-time}, the sum in
\eqref{eq:block-witness-sum} equals the intrinsic measurable sum
\begin{equation}\label{eq:intrinsic-moving-block-sum}
 \sum_{\ell'\in B_{S_0^{-r}b}(u)}
 \left|\int f(T_X^r x(y))
                  \overline{\chi_{u,\ell'}(y)}\dd\eta_u(y)\right|.
\end{equation}
Only the set of lines is relevant; changing representative phases
or permuting labels inside this set changes no summand magnitude.
This proves measurability without a measurable choice of $J_u$.

The conditional laws are equivariant:
$(T_Y^r)_*\eta_u=\eta_{T_U^r u}$.
Time carries $B_{S_0^{-r}b}(u)$ bijectively onto
$B_b(T_U^r u)$, taking representatives to representatives up to
scalar phase. The change of variables $y'=T_Y^r y$ in
\eqref{eq:intrinsic-moving-block-sum} consequently gives
\begin{equation}\label{eq:witness-block-transport}
 s_r(u)=c_b(T_U^r u).
\end{equation}
Invariance of $\nu$ proves its integral is $c$ for every $r$.
Finally, $G^r$ bijects $B$ onto $B_{S_0^{-r}b}$, so every moving
block has the same cardinality $|B|=m_b$. Each coefficient is at
most $\norm f_\infty$. This proves the bound in
\eqref{eq:block-witness-integral} and completes the Proposition.
\end{proof}

Proposition~\ref{prop:block-witness} leaves a problem on a single
probability fibre: show that the absolute coefficients along each
orbit $G^r\ell$ have Ces\`aro mean zero. The uniform finite-block
bound will then allow integration over $u$. The next section proves
this cancellation using only the mixing zero-entropy $X$ marginal
and the group relations \eqref{eq:block-character-laws}.
